\documentclass[11pt]{article}
\usepackage{marcoreport}
\usepackage{listings}
\usepackage{xcolor}

% Code listing style (inherited from tool-exploration)
\lstset{
  language=Python,
  basicstyle=\ttfamily\small,
  keywordstyle=\color{blue!70!black},
  commentstyle=\color{green!50!black},
  stringstyle=\color{red!60!black},
  backgroundcolor=\color{gray!5},
  frame=single,
  framerule=0.4pt,
  rulecolor=\color{gray!50},
  breaklines=true,
  showstringspaces=false,
  tabsize=4,
  aboveskip=0.8\baselineskip,
  belowskip=0.8\baselineskip
}

\title{Coverage Validation: \texttt{edgar-crawler} vs Hoberg-Phillips TNIC}
\author{Marco Zhang}
\date{February 2026}

\begin{document}
\maketitle

% ---------------------------------------------------------------------------
% Section 1: Where We Are & Goal
% ---------------------------------------------------------------------------
\section{Current Stage \& Goal}

In the previous report, we explored \texttt{edgar-crawler}'s mechanics: how it
downloads filings from EDGAR, parses raw HTML/txt, and extracts item-level text.
We documented configuration caveats (exact filing-type matching misses 47\% of
variants, item lists need manual upkeep across SEC eras) and confirmed extraction
reliability across 1993--2024.

The next question: \textbf{how does edgar-crawler's coverage compare to an
established benchmark?} We choose Hoberg-Phillips TNIC (2005) because (1)~it is
widely used in finance/accounting research, (2)~it is built from 10-K Item~1
extractions---the same task \texttt{edgar-crawler} performs, and (3)~it publishes
both data (pairwise similarity scores by gvkey-year) and replication code (Python
notebooks). This lets us validate at two levels: file-level (do both tools see
the same firms?) and item-level (do both tools extract the same text?).

% ---------------------------------------------------------------------------
% Section 2: File-Level Coverage
% ---------------------------------------------------------------------------
\section{File-Level Coverage}

TNIC uses Compustat gvkeys---its universe is public companies traded on the
three major US exchanges. EDGAR uses CIK as identifier, covering all SEC
registrants (small business filers, foreign issuers, shell companies), so it is
much broader.

We match EDGAR filings to TNIC firms using the WRDS \texttt{wciklink\_gvkey}
crosswalk (historical CIK-gvkey mapping from the SEC Analytics Suite). We filter
EDGAR filings by \texttt{year(datadate) == 2005}, which exactly matches TNIC's
year definition---confirmed from H-P's own
README\footnote{\url{https://hobergphillips.tuck.dartmouth.edu/idata/Readme_tnic3HHIData.txt}}:
\emph{``the year field in this database is based on Compustat calendar years
obtained as the first four digits of the YYYYMMDD datadate variable.''}

\begin{table}[h]
\centering
\begin{tabular}{lrrl}
\toprule
Method & FY2005 CIKs & TNIC Covered & Rate \\
\midrule
SEC API + pagination & 12,038 & 5,084 / 5,122 & \textbf{99.3\%} \\
\bottomrule
\end{tabular}
\caption{File-level coverage: EDGAR vs TNIC (2005).}
\end{table}

This is a strict, exact-fiscal-period match---not a loose filing-window
heuristic---making 99.3\% a reliable statistic. Good enough to proceed from
file-level to item-level comparison.

% --- Filing type breakdown ---
\subsection{Filing Type Breakdown}

Among the 5,137 TNIC-matched CIKs with FY2005 filings, \textbf{$\sim$94\% file
10-K and $\sim$6\% file 10KSB}:

\begin{table}[h]
\centering
\begin{tabular}{lrl}
\toprule
Filing Type & TNIC-Matched Firms & \% \\
\midrule
10-K / 10-K/A   & 4,824 & 93.9\% \\
10KSB / 10KSB/A & 313   & 6.1\%  \\
\bottomrule
\end{tabular}
\caption{Filing type distribution among TNIC-matched FY2005 firms (Compustat universe).}
\end{table}

This tells us that for pre-2009 data (before the SEC eliminated the 10KSB form),
we need parsing methodology that handles 10KSB filings---not just 10-K.

If we want to use the \textbf{full EDGAR CIK universe} (not just TNIC-matched),
the 10KSB share is even larger:

\begin{table}[h]
\centering
\begin{tabular}{lrl}
\toprule
Filing Type & Count & Share \\
\midrule
10-K      & 9,003 & 59.9\% \\
10KSB     & 3,223 & 21.5\% \\
10-K/A    & 1,517 & 10.1\% \\
10KSB/A   & 1,262 & 8.4\% \\
10-KT / 10-KT/A & 19 & 0.1\% \\
\midrule
\textbf{10-K variants} & \textbf{10,520} & \textbf{73.7\%} \\
\textbf{10KSB variants} & \textbf{4,485} & \textbf{26.2\%} \\
\bottomrule
\end{tabular}
\caption{FY2005 filing type breakdown across all EDGAR CIKs.}
\end{table}

So for full-universe coverage, we need parsing methodology for more than just
10-K and 10KSB---amendment variants (10-K/A, 10KSB/A) and transition reports
(10-KT) also need handling.

% --- The last 0.7% ---
\subsection{The Last 0.7\%}

We diagnosed all 38 NOT\_COVERED gvkeys. The largest category (17/38 = 45\%) is
\textbf{SEC API data quality issues}---the filings exist in EDGAR, but the SEC
API returns an incorrect \texttt{reportDate} (equal to the filing date, or off by
one day), so our year-matching filter misses them.

\begin{table}[h]
\centering
\small
\begin{tabular}{lrll}
\toprule
Category & Count & Key Insight & Example \\
\midrule
SEC rd $\approx$ filingDate & 14 & Dec FY-end; SEC returns rd = fd & Ball Corp \\
SEC rd off-by-one           & 3  & rd = 2006-01-01 not 2005-12-31 & Interface Inc \\
Genuine non-Dec FY          & 6  & 52/53-wk years or Jan/Mar FY   & Yum Brands \\
Missing FY2005 filing       & 5  & No 10-K for FY2005 in EDGAR    & Nature's Sunshine \\
Late filer                  & 4  & FY05 10-K filed years late      & Am.\ Italian Pasta \\
Spinoff / new entity        & 4  & CIK didn't exist until 2006+   & Celera Corp \\
Foreign private issuer      & 1  & Files 6-K only, not 10-K       & AXA S.A. \\
No crosswalk                & 1  & No wciklink entry              & gvkey 62642 \\
\midrule
\textbf{Total}              & \textbf{38} & & \\
\bottomrule
\end{tabular}
\caption{Diagnosis of all 38 NOT\_COVERED gvkeys (0.7\% gap).}
\end{table}

% ---------------------------------------------------------------------------
% Section 3: Item-Level Extraction
% ---------------------------------------------------------------------------
\section{Item-Level Extraction}

H-P only extract Item~1 (Business Description) and only provide replication code
for that item, so we focus our comparison on Item~1. We use
\texttt{edgar-crawler}'s 62 test fixture 10-K filings (spanning 1993--2023, mixed
\texttt{.txt}/\texttt{.htm} formats) and run both parsers on each filing:

\begin{itemize}
  \item \textbf{edgar-crawler}: \texttt{ExtractItems} class with
    \texttt{items\_to\_extract=["1"]}, \texttt{lxml} parser, newline-anchored
    regex.
  \item \textbf{H-P}: reimplemented from their Notebook~2---4-regex cascade,
    BeautifulSoup \texttt{html.parser}, longest-match disambiguation, 2KB minimum
    filter.
\end{itemize}

We compare the extracted Item~1 text using word-level Jaccard similarity.

% --- Results ---
\subsection{Results}

\textbf{Extraction success rate}:
\begin{itemize}
  \item edgar-crawler: \textbf{62/62 (100\%)}
  \item H-P: \textbf{35/62 (56.5\%)}
\end{itemize}

\textbf{Jaccard similarity} (35 filings where both extracted):
Mean = 0.910, Median = 0.946, Min = 0.241, Max = 0.999.

\begin{table}[h]
\centering
\begin{tabular}{lrll}
\toprule
Category & Count & Share & Description \\
\midrule
MATCH          & 16 & 16/35 & Jaccard $\geq$ 0.95, essentially identical \\
BOUNDARY\_DIFF & 17 & 17/35 & Jaccard 0.70--0.95, same core, different boundaries \\
MAJOR\_DIFF    & 2  & 2/35  & Jaccard $<$ 0.70, substantially different \\
EC\_ONLY       & 27 & ---   & edgar-crawler extracted, H-P did not \\
HP\_ONLY       & 0  & ---   & H-P extracted, edgar-crawler did not \\
BOTH\_EMPTY    & 0  & ---   & Neither extracted anything \\
\bottomrule
\end{tabular}
\caption{Item~1 extraction comparison across 62 test fixture filings.}
\end{table}

% --- Key Findings ---
\subsection{Key Findings}

\textbf{1. H-P fails on all pre-2004 filings (27 EC\_ONLY).}
H-P's regex looks for \texttt{item 1 ... item 1a/1b} boundaries, but many older
filings go directly from ``ITEM~1. BUSINESS'' to ``ITEM~2. PROPERTIES''---no
Item~1A ``Risk Factors'' (not mandatory until the 2005 SEC rule change). H-P
cannot find an ending boundary, so it returns nothing. edgar-crawler handles this
by trying all subsequent items as end boundaries
(1A $\to$ 1B $\to$ 2 $\to$ 3 $\to$ \ldots $\to$ SIGNATURE).

\textbf{2. H-P includes Table of Contents pollution (BOUNDARY\_DIFF).}
In 15 of 17 BOUNDARY\_DIFF cases, H-P's extraction starts earlier---capturing
``Item~1'' from the Table of Contents and extending to ``Item~1A'' in the body.
edgar-crawler avoids this via newline-anchored regex (\texttt{\textbackslash n}
at start of item header), which filters out inline ToC references. For modern
iXBRL filings (2019+), H-P also captures raw HTML/CSS because
\texttt{BeautifulSoup('html.parser')} doesn't fully parse inline XBRL tags.

\textbf{3. Two MAJOR\_DIFF cases.}
\begin{itemize}
  \item Horizon Financial 2009 (J=0.24): H-P extracted 15KB (ToC fragment only),
    edgar-crawler extracted 139KB (full Item~1).
  \item FedEx 2023 (J=0.34): H-P extracted 530KB (raw HTML including
    CSS/iXBRL), edgar-crawler extracted 112KB (clean text).
\end{itemize}

\textbf{4. When both succeed, they agree closely.}
For the 33 filings with Jaccard $\geq$ 0.70, the mean Jaccard is 0.946.
Remaining differences are boundary choices (H-P includes the trailing ``Item~1A''
header) and table removal (edgar-crawler strips numerical tables).

\textbf{5. Year distribution confirms the pattern is era-driven.}

\begin{table}[h]
\centering
\begin{tabular}{lrrrrr}
\toprule
Era & MATCH & BOUNDARY\_DIFF & MAJOR\_DIFF & EC\_ONLY & Total \\
\midrule
1993--1999 & 0  & 0  & 0 & 14 & 14 \\
2000--2004 & 1  & 0  & 0 & 9  & 10 \\
2005--2010 & 3  & 5  & 1 & 3  & 12 \\
2011--2023 & 12 & 12 & 1 & 1  & 26 \\
\bottomrule
\end{tabular}
\caption{Extraction comparison by era. Pre-2005: H-P extracts almost nothing
(23/24 EC\_ONLY). Post-2010: both work on 24/26 filings (92\%).}
\end{table}

Pre-2005: H-P extracts almost nothing---no Item~1A boundary available. 2005--2010:
transition period, H-P starts working. 2011--2023: both work on 24/26 filings
(92\%); BOUNDARY\_DIFFs are consistently ToC pollution. The 2005 cutoff aligns
with SEC's mandatory Item~1A (Risk Factors) requirement.

\end{document}
